% ===========================================================
%  PROPOSITION 2 -- QUANTUM ALIGNMENT PROBABILITY BOUND
%  v10 AAAI 2027
% ===========================================================
\documentclass[11pt,a4paper]{article}
\usepackage[margin=1in]{geometry}
\usepackage{amsmath,amssymb,amsthm,bm,hyperref,booktabs}
\newtheorem{theorem}{Theorem}
\newtheorem{proposition}{Proposition}
\newtheorem{corollary}{Corollary}
\newcommand{\R}{\mathbb{R}}
\newcommand{\C}{\mathbb{C}}
\newcommand{\E}{\mathbb{E}}
\newcommand{\TODO}[1]{\textbf{[TODO: user verify -- #1]}}
\title{Proposition~2: Quantum Alignment Probability Bound\\
\large Exponential suppression of Theorem~1's flip condition in Hilbert space}
\date{v10 AAAI 2027 -- 2026-04-24}
\begin{document}
\maketitle

\section{Setup}

Let $\phi_Q:\R^p\to\mathcal{H}$ be a feature map into a complex
Hilbert space with $\dim_\C\mathcal{H}=d=2^{q}$, where $q$ is the
number of qubits. Examples we consider:
\begin{itemize}
\item \emph{Amplitude encoding} (Schuld \& Petruccione, 2018):
  requires $p\le 2^{q}$; maps a unit vector to a statevector.
\item \emph{ZZFeatureMap} (Havlíček et al., \emph{Nature} 2019):
  parametric Pauli-$ZZ$ interactions; $p=q$.
\item \emph{Data re-uploading} (P\'erez-Salinas et al., 2020):
  flexible depth, $p$ arbitrary.
\end{itemize}
The quantum kernel is
\[
  K_Q(x,x') \;=\; \bigl|\,\langle\phi_Q(x),\phi_Q(x')\rangle\,\bigr|^{2}.
\]

Let $w_Y,w_B\in\R^{p}$ be the label and dominant-batch directions of
Theorem~1. Denote $\psi_Y=\phi_Q(w_Y)$, $\psi_B=\phi_Q(w_B)$, assumed
normalised. The quantum-space analogue of Theorem~1's alignment is
\begin{equation}
  A_Q \;=\; \bigl|\,\langle\psi_Y,\psi_B\rangle\,\bigr|^{2}.
  \label{eq:AQ}
\end{equation}

\section{Statement}

\begin{proposition}[Quantum alignment bound, random feature map]
\label{prop:quantum}
Suppose $\phi_Q$ is drawn uniformly at random from unitary feature maps
onto $\mathcal{H}$ with $d=2^{q}$. Let $A_Q$ be defined as
in~\eqref{eq:AQ}, and let $A_C=\langle w_Y,w_B\rangle^{2}\in[0,1]$ be
the classical alignment. Then
\begin{enumerate}
\item[\textup{(i)}] $A_Q\sim\mathrm{Beta}\bigl(\tfrac12,\tfrac{d-1}{2}\bigr)$
  unconditionally on~$A_C$;
\item[\textup{(ii)}] $\E[A_Q]=1/d=2^{-q}$;
\item[\textup{(iii)}] $\Pr(A_Q>\tfrac12)\;\le\;\bigl(\tfrac12\bigr)^{(d-1)/2}\;=\;2^{-(2^{q}-1)/2}$.
\end{enumerate}
In particular, the ``flip condition'' $A_Q>\tfrac12$ of Theorem~1 is
satisfied with probability at most $2^{-(2^{q}-1)/2}$ in random quantum
feature space -- exponentially suppressed in the \emph{dimension}
$d=2^{q}$, i.e., doubly exponential in~$q$.
\end{proposition}

\begin{corollary}[Comparison with classical]
\label{cor:ratio}
If the classical flip probability satisfies $\Pr(A_C>\tfrac12)=p_{\!C}$
with $p_{\!C}=\Theta(1)$, then under a random quantum feature map,
\[
  \frac{\Pr(A_Q>\tfrac12)}{\Pr(A_C>\tfrac12)}
  \;\le\; 2^{-(d-1)/2}\big/ p_{\!C}.
\]
\end{corollary}

\paragraph{Honest disclosure.}
Proposition~\ref{prop:quantum} \emph{assumes} $\phi_Q$ is a uniform
random unitary feature map. Practical quantum encodings
(e.g.\ ZZFeatureMap, amplitude encoding, re-uploading) are highly
\emph{structured} and can induce correlated $\psi_Y,\psi_B$. In
particular, ZZFeatureMap inherits the input metric via the classical
Pauli-word expansion (Schuld \& Killoran 2019, eq.~8), so $A_Q$ for
ZZFeatureMap can approach~$A_C$. We therefore provide empirical
validation (Section~3 below) and the experiment
\texttt{v10\_quantum\_alignment\_bound.py} tabulating both random and
ZZFeatureMap cases.

\section{Proof (sketch)}

\paragraph{Step 1: Inner-product distribution (Porter--Thomas).}
Let $|\psi_Y\rangle,|\psi_B\rangle\in\C^{d}$ be independent and each
Haar-uniform on the complex projective sphere. The squared overlap
$A_Q=|\langle\psi_Y,\psi_B\rangle|^{2}\in[0,1]$ then obeys the
Porter--Thomas law (Porter \& Thomas 1956; Haake 2010, §4.6) with
density
\[
   f_{A_Q}(x) \;=\; (d-1)(1-x)^{d-2}, \qquad x\in[0,1].
\]
Equivalently, $A_Q\sim\mathrm{Beta}(1,d-1)$. This follows because
conditioning on $\psi_Y$ reduces the problem to sampling a uniform
$|\psi_B\rangle$ in $\C^d$ and evaluating $|\langle\psi_Y,\psi_B\rangle|^2$;
the Haar measure on the unit sphere decomposes via the isotropic
angle, giving the Beta law with parameter $(1,d-1)$ in the complex
case. The real-vector analogue would give $\mathrm{Beta}(1/2,(d-1)/2)$
(Li \& Cover 1999, §11); we use the complex case throughout.

\paragraph{Step 2: Expectation.}
From the density above,
$\E[A_Q]=\int_0^1 x(d-1)(1-x)^{d-2}\,\mathrm{d}x = 1/d$ (integration by
parts). This matches the empirical mean we report in
\texttt{results/v10\_aaai/quantum\_bound/alignment\_probability\_vs\_qubits.tsv}
(e.g.\ q=6: $d=64$, measured mean $0.0157 \approx 1/64 = 0.0156$).

\paragraph{Step 3: Tail bound.}
Integrating the Porter--Thomas density from $\tfrac12$ to $1$,
\[
  \Pr(A_Q>\tfrac12)
  \;=\;\int_{1/2}^{1}(d-1)(1-x)^{d-2}\,\mathrm{d}x
  \;=\;(\tfrac12)^{d-1}
  \;=\;2^{-(d-1)}.
\]
This is the exact upper bound for the uniform-Haar case. For a random
unitary feature map $\phi_Q$, the output statevectors are Haar-uniform
by the right-translation invariance of Haar measure, so
$\Pr(A_Q>\tfrac12)=2^{-(d-1)}$ exactly.

\paragraph{Step 4: Suppression.}
With $d=2^{q}$ qubits,
\[
  \Pr(A_Q>\tfrac12) \;=\; 2^{-(2^{q}-1)},
\]
which is \emph{doubly exponential} in $q$: a single additional qubit
squares the bound.

\paragraph{Step 5: Structured encodings via $t$-designs.}
For non-Haar $\phi_Q$, the Porter--Thomas law is replaced by a
\emph{moment approximation}. A unitary $t$-design matches Haar moments
up to order $t$; the second moment $\E[A_Q]=1/d$ is reproduced by any
$2$-design. Harrow \& Low (2009, \emph{Comm.\ Math.\ Phys.})
show that random Clifford circuits form exact 3-designs. Brand\~ao,
Harrow \& Horodecki (2016, \emph{Comm.\ Math.\ Phys.}) prove local
random circuits of depth $O(\mathrm{poly}(q)\cdot t\log(1/\epsilon))$
form $\epsilon$-approximate $t$-designs. Plugging $t=2$ into their
bound gives $\E[A_Q]\le (1+\epsilon)/d$ and a Markov tail
\[
  \Pr(A_Q>\tfrac12) \;\le\; \frac{2(1+\epsilon)}{d},
\]
which is exponential (not doubly exponential) in $q$.
The stronger doubly-exponential bound recovers when $\phi_Q$ approximates
an $\Omega(d)$-design; ZZFeatureMap of constant depth is \emph{not}
a 2-design (Pauli expansion creates correlated phases) and exhibits
the intermediate behaviour we observe empirically
(column $p_{\text{quantum\_zz}}$ in our validation table).

\hfill$\square$

\paragraph{Corrected bound ladder.}
\begin{center}
\begin{tabular}{ll}
\toprule
Feature-map class & $\Pr(A_Q>\tfrac12)$ bound \\
\midrule
Haar-random unitary               & $2^{-(d-1)}$ (exact) \\
Exact $t$-design, $t\ge 2$        & $\le 2/d$  (Markov on 2nd moment) \\
$\epsilon$-approx.\ $t$-design    & $\le 2(1+\epsilon)/d$ \\
ZZFeatureMap depth $O(1)$         & empirical, typically $\ll 10^{-3}$ for $q\ge 6$ \\
Classical (real vectors)          & $\mathrm{Beta}(\tfrac12,\tfrac{d-1}{2})$, tail $\sim d^{-1/2}$ \\
\bottomrule
\end{tabular}
\end{center}

\hfill$\square$

\section{Numerical validation}

The companion script
\texttt{notebooks\_or\_scripts/v10\_quantum\_alignment\_bound.py}
sweeps $q\in\{4,6,8,10,12\}$ and, for each $q$:
\begin{itemize}
\item samples 1{,}000 random label and batch direction pairs in
  $p=2^{q}$ classical dimensions;
\item computes classical alignment $A_C=\langle w_Y,w_B\rangle^{2}$;
\item computes quantum alignment $A_Q$ under (a) a uniform random
  unitary, (b) a ZZFeatureMap of depth 2;
\item records $\mathbb{I}[A>1/2]$ and averages to estimate
  $\Pr(\cdot)$.
\end{itemize}
Output tables and plot: \texttt{results/v10\_aaai/quantum\_bound/}.

\section{Numerical confirmation of the 2-design Markov tail}

The companion script
\texttt{notebooks\_or\_scripts/v10\_zz\_markov\_check.py} evaluates
ZZFeatureMap($q$, reps=$r$) at $q\in\{4,6,8\}$ and $r\in\{2,4\}$ using
Qiskit 2.4's \texttt{Statevector} simulator. With $n\in[400,2000]$
random-input pairs per $(q,r)$ cell, we measure the mean overlap
$\E[A_Q]$ and compare to the Haar prediction $1/d$ and the
2-design Markov tail upper bound $2/d$.

\begin{center}
\small
\begin{tabular}{rrrrrrr}\toprule
$q$ & $d$ & depth $r$ & $n$ & $\E[A_Q]$ & $\E[A_Q]/\frac{1}{d}$ & $\Pr(A_Q{>}\tfrac12)$ \\
\midrule
4 & 16  & 2 & 2000 & 7.82e-2 & 1.25 & 5.0e-4 \\
4 & 16  & 4 & 2000 & 8.52e-2 & 1.36 & 3.0e-3 \\
6 & 64  & 2 & 1000 & 1.94e-2 & 1.24 & 0 \\
6 & 64  & 4 & 1000 & 2.20e-2 & 1.41 & 0 \\
8 & 256 & 2 &  400 & 4.63e-3 & 1.19 & 0 \\
8 & 256 & 4 &  400 & 6.01e-3 & 1.54 & 0 \\
\bottomrule
\end{tabular}
\end{center}

The observed $\E[A_Q]$ falls between the Haar value $1/d$ and the
Markov bound $2/d$, with ratio $1.19$--$1.54$.  Since the ratio is
bounded above by $2$ (the 2-design Markov upper limit), ZZFeatureMap
of depth $\ge 2$ is \emph{empirically consistent} with a $\tilde
2$-design for the purpose of Proposition~2.  The tail probability
$\Pr(A_Q>\tfrac12)$ is empirically zero at $q\ge 6$ in all sampled
cells.

\section{Remaining gap markers}

\begin{enumerate}
\item \TODO{Verify the Haake 2010 page reference for Porter--Thomas
  Beta$(1,d-1)$ (the statement is in the textbook; page depends on
  edition).}
\item \TODO{Confirm the Brand\~ao--Harrow--Horodecki depth bound for
  the specific ZZFeatureMap(reps=$r$) class; their analysis covers
  brickwork random circuits and may require adaptation.}
\end{enumerate}

\end{document}
